
% --- Párrafo 1: qué son los TEs --------------------------------------
Los elementos transponibles (TEs, por sus siglas en inglés) son secuencias
de ADN capaces de moverse o de copiarse dentro del genoma que las
hospeda~\cite{bourque2018ten}. Constituyen una fracción considerable de
muchos genomas eucariotas: cerca de la mitad del genoma humano deriva de
ellos~\cite{lander2001initial} y en el maíz representan alrededor del
85\,\% de la secuencia~\cite{schnable2009b73}. Lejos de ser ``ADN
basura'', los TEs participan en la evolución de los genomas, en la
regulación de la expresión génica y en la generación de variación
genética sobre la que actúa la selección~\cite{bourque2018ten}.
% Fuente: Bourque et al. 2018 (definición, papel evolutivo y regulatorio);
%         Lander et al. 2001 (~45 % del genoma humano);
%         Schnable et al. 2009 (~85 % del genoma del maíz B73).

% --- Párrafo 2: cómo se clasifican -----------------------------------
Para estudiarlos, los TEs se organizan en una clasificación jerárquica.
El sistema propuesto por Wicker et al.~\cite{wicker2007unified} los
divide en dos clases según su mecanismo de transposición: la Clase~I,
o retrotransposones, que se movilizan a través de un intermediario de
ARN, y la Clase~II, o transposones de ADN. Cada clase se subdivide en
órdenes (por ejemplo, LTR, LINE o TIR) y estos, a su vez, en
superfamilias (por ejemplo, \textit{Gypsy}, \textit{Copia} o
\textit{hAT}). Esta jerarquía Clase/Orden/Superfamilia es la que siguen
las etiquetas de la base de datos utilizada en este
trabajo~\cite{orozcoarias2026panteon}.
% Fuente: Wicker et al. 2007 (Clase I/II, orden, superfamilia);
%         PanTEon (bioRxiv) usa esa jerarquía de 3 niveles.

% --- Párrafo 3: por qué es un reto almacenarlos y analizarlos ---------
El análisis de los TEs plantea retos tanto biológicos como
computacionales. Desde el punto de vista biológico, los TEs suelen ser
específicos de cada especie, divergen con rapidez incluso entre especies
cercanas y, dentro de una misma familia, sus copias van desde casi
idénticas hasta muy degradadas; además, su alto número de copias y las
inserciones anidadas dificultan su detección y anotación~\cite{orozcoarias2026panteon}.
Desde el punto de vista de los datos, la información sobre TEs es
voluminosa, heterogénea y poco estructurada: las secuencias se
distribuyen en archivos FASTA cuyo encabezado mezcla en texto libre el
identificador, la clasificación y la especie, mientras que la taxonomía
de cada especie y el origen de cada biblioteca se describen en archivos
de metadatos aparte~\cite{orozcoarias2026panteon,panteon2026zenodo}.
A ello se suma la transformación habitual de las secuencias en perfiles
de \textit{k}-mers, una de las representaciones más usadas para
entrenar clasificadores de TEs~\cite{orozcoarias2026panteon}: con el
alfabeto \{A, C, G, T\} existen $4^k$ \textit{k}-mers posibles, de modo
que el perfil completo para $k=1,\dots,6$ tiene hasta 5\,460 valores por
secuencia. Para las 239\,946 secuencias de PanTEon v1.6.1 esto supone
del orden de $1{,}3\times10^{9}$ valores, lo que convierte la forma de
almacenarlos y consultarlos en una decisión de diseño central.
% Fuente: PanTEon (bioRxiv) – especificidad, divergencia, copias,
%         anidamiento, datos categóricos y poco estructurados, k-mers
%         usados por 5 de 7 herramientas revisadas.
% Cálculo propio: 4+16+64+256+1024+4096 = 5460;
%                 239946 × 5460 = 1 310 105 160 ≈ 1,3 × 10^9.
% Conteo propio: 239 946 secuencias (grep -c "^>" sobre el FASTA).

% --- Párrafo 4: objetivo del trabajo ----------------------------------
En este trabajo diseñamos, implementamos y evaluamos una base de datos
documental en MongoDB que integra las secuencias de TEs de PanTEon
Database v1.6.1~\cite{panteon2026zenodo}, la taxonomía de las especies
de las que proceden y sus perfiles de \textit{k}-mers ($k=1$ a $6$).
El procesamiento se realiza con un \textit{pipeline} reproducible en
Python que lee y valida las secuencias, documenta y resuelve los
problemas de calidad de las fuentes, cruza ambas fuentes por especie y
carga los datos de forma masiva. Además, respondemos un conjunto de
preguntas de investigación mediante consultas y agregaciones, y
comparamos experimentalmente dos estrategias de almacenamiento de los
\textit{k}-mers en términos de tiempo de carga, espacio en disco y
tiempo de consulta.

% =====================================================================
\subsection*{Trabajo relacionado}

% --- PanTEon -----------------------------------------------------------
\paragraph{PanTEon.}
PanTEon~\cite{orozcoarias2026panteon} es un marco de trabajo para
diseñar y comparar clasificadores de TEs basados en aprendizaje
profundo. Combina una base de datos de TEs curada automáticamente, que
abarca los reinos Animalia, Plantae y Fungi, con una plataforma que
estandariza el entrenamiento, la evaluación y la inferencia de nueve
métodos de aprendizaje automático. Su versión 1.6.1 reúne 239\,946 TEs
de 2\,790 especies. Las secuencias proceden de dos orígenes: familias
de Dfam y bibliotecas generadas \textit{de novo} con
RepeatModeler2~\cite{flynn2020repeatmodeler2} y curadas
automáticamente con MCHelper~\cite{orozcoarias2024mchelper}; solo se
conservaron las secuencias con todos los rasgos estructurales esperados
para su superfamilia. Cada secuencia se etiqueta con la jerarquía
Clase/Orden/Superfamilia. Los datos se distribuyen en dos archivos
planos: un FASTA con las secuencias y un CSV con la taxonomía y el
origen de cada biblioteca, relacionados únicamente por el nombre de la
especie~\cite{panteon2026zenodo}. Esta organización es adecuada para
entrenar modelos, pero no ofrece consultas, índices ni relaciones
explícitas entre secuencias y taxonomía, y no almacena los perfiles de
\textit{k}-mers, que cada usuario debe recalcular.
% Fuente: PanTEon (bioRxiv) – marco, 3 reinos, 9 métodos, 239 946 TEs,
%         2 790 especies, Dfam + RM2 + MCHelper, filtros estructurales,
%         jerarquía de 3 niveles.
% Fuente: Zenodo + nuestra exploración – 2 archivos planos unidos por
%         el nombre de la especie (CSV: 2 404 especies RM2+MCH y 386 Dfam).
% La última frase es la brecha que justifica nuestro trabajo.

% --- Dfam --------------------------------------------------------------
\paragraph{Dfam.}
Dfam~\cite{storer2021dfam} es un recurso comunitario de familias de TEs
y otros repetidos. Cada familia se construye a partir de un alineamiento
semilla (\textit{seed alignment}) de copias representativas, del que se
derivan un modelo oculto de Markov de perfil (HMM) y una secuencia
consenso; ambos se usan para anotar genomas. Las familias se dividen en
curadas, con umbrales de tasa de falsos positivos caracterizados, y no
curadas, identificadas con el prefijo de acceso ``DR''. Cada familia se
asocia a la taxonomía del NCBI, y la clasificación se representa como
un árbol no jerarquizado por rangos cuyos nombres se concatenan desde la
raíz (por ejemplo, \texttt{Interspersed\_repeat;Unknown}), con
correspondencias hacia los esquemas de RepeatMasker, Repbase y Wicker
et al. Dfam se consulta mediante un portal web y una API, y se
distribuye también en FamDB, un formato basado en HDF5 que permite
consultas locales. En la versión descrita por Storer et
al.~\cite{storer2021dfam} contenía 273\,655 familias.
% Fuente: Storer et al. 2021 (Mobile DNA) – seed alignment, HMM y
%         consenso, curadas vs. no curadas (prefijo DR), taxonomía NCBI,
%         árbol de clasificación no jerarquizado, mapeos, API, FamDB
%         (HDF5), 273 655 familias (versión 3.3).

% --- Contraste y aporte ------------------------------------------------
\paragraph{Diferencia con este trabajo.}
Dfam está orientado a la anotación de genomas y representa cada familia
mediante modelos (HMM y consenso), mientras que PanTEon está orientado
al entrenamiento de clasificadores y distribuye las secuencias en
archivos planos. Ninguno de los dos almacena perfiles de \textit{k}-mers
consultables junto con la secuencia y la taxonomía de la especie. Este
trabajo aborda esa necesidad mediante un modelo documental que permite
recorrer la relación secuencia--especie en ambos sentidos, consultar
los \textit{k}-mers de cualquier secuencia y evaluar el costo de
distintas formas de almacenarlos.